\documentclass[12pt,a4paper]{article}

\usepackage[utf8]{inputenc}
\usepackage[T1]{fontenc}
\usepackage{geometry}
\usepackage{hyperref}
\usepackage{listings}
\usepackage{xcolor}
\usepackage{titlesec}
\usepackage{fancyhdr}
\usepackage{graphicx}
\usepackage{booktabs}
\usepackage{enumitem}
\usepackage{parskip}
\usepackage{amsmath}
\usepackage{tocloft}
\usepackage{mdframed}

\geometry{top=2.5cm, bottom=2.5cm, left=2.5cm, right=2.5cm}

% Colors
\definecolor{codegreen}{rgb}{0,0.6,0}
\definecolor{codegray}{rgb}{0.5,0.5,0.5}
\definecolor{codepurple}{rgb}{0.58,0,0.82}
\definecolor{backcolour}{rgb}{0.97,0.97,0.97}
\definecolor{accentblue}{RGB}{30,90,160}
\definecolor{lightblue}{RGB}{220,235,250}

% Code listing style
\lstdefinestyle{mystyle}{
    backgroundcolor=\color{backcolour},
    commentstyle=\color{codegreen},
    keywordstyle=\color{accentblue}\bfseries,
    numberstyle=\tiny\color{codegray},
    stringstyle=\color{codepurple},
    basicstyle=\ttfamily\footnotesize,
    breakatwhitespace=false,
    breaklines=true,
    captionpos=b,
    keepspaces=true,
    numbers=left,
    numbersep=5pt,
    showspaces=false,
    showstringspaces=false,
    showtabs=false,
    tabsize=2,
    frame=single,
    rulecolor=\color{accentblue!40}
}
\lstset{style=mystyle}

% Section formatting
\titleformat{\section}{\large\bfseries\color{accentblue}}{{\thesection}}{1em}{}[\color{accentblue!40}\titlerule]
\titleformat{\subsection}{\normalsize\bfseries\color{accentblue!80}}{\thesubsection}{1em}{}
\titleformat{\subsubsection}{\normalsize\bfseries\color{accentblue!60}}{\thesubsubsection}{1em}{}

% Header/Footer
\pagestyle{fancy}
\fancyhf{}
\fancyhead[L]{\small\textcolor{accentblue}{Mini Game Hub --- SSL Project 2026}}
\fancyhead[R]{\small\textcolor{accentblue}{Poladi \& Singh}}
\fancyfoot[C]{\small\thepage}
\renewcommand{\headrulewidth}{0.4pt}
\renewcommand{\footrulewidth}{0.4pt}
\setlength{\headheight}{15pt}

\hypersetup{
    colorlinks=true,
    linkcolor=accentblue,
    urlcolor=accentblue,
    citecolor=accentblue,
    pdftitle={Mini Game Hub - SSL Project Report},
    pdfauthor={Samhitha Poladi, Riddhima Singh}
}

% -----------------------------------------------
\begin{document}

% Title Page
\begin{titlepage}
    \centering
    \vspace*{2cm}
    {\Huge\bfseries\color{accentblue} Mini Game Hub\par}
    \vspace{0.5cm}
    {\Large SSL Project Report --- 2026\par}
    \vspace{2cm}
    \begin{mdframed}[backgroundcolor=lightblue, linecolor=accentblue, linewidth=1.5pt]
        \centering
        \vspace{0.3cm}
        {\large\bfseries Submitted by}\par
        \vspace{0.4cm}
        {\large Samhitha Poladi \quad (25B1088)}\par
        \vspace{0.2cm}
        {\large Riddhima Singh \quad\enspace (25B1068)}\par
        \vspace{0.3cm}
    \end{mdframed}
    \vspace{2cm}
    {\large\bfseries Abstract}\par
    \vspace{0.5cm}
    \begin{minipage}{0.85\textwidth}
        \normalsize
        Mini Game Hub is a two-player interactive gaming platform that integrates Bash scripting
        with Python-based graphical games. The system provides secure user authentication, a
        selection of classic board games rendered via Pygame, persistent game history, a
        statistical leaderboard, and data visualisations. This report documents the external
        libraries used, features implemented, code logic for each component, challenges
        encountered, and reflections on future improvements.
    \end{minipage}
    \vfill
    {\large April 2026\par}
\end{titlepage}

\tableofcontents
\newpage

% -----------------------------------------------
\section{External Libraries Used and Their Purpose}

The project uses several well-established external libraries, each serving a specific role in the system.

\begin{center}
\begin{tabular}{lll}
\toprule
\textbf{Library} & \textbf{Language} & \textbf{Purpose} \\
\midrule
\texttt{pygame}       & Python & GUI rendering, event handling, game windows \\
\texttt{numpy}        & Python & Board representation, win-checking logic \\
\texttt{matplotlib}   & Python & Data visualisations (pie chart, bar chart) \\
\texttt{hashlib / shasum} & Python / Bash & Password hashing (SHA-256) \\
\texttt{os}           & Python & File path management across platforms \\
\texttt{time}         & Python & End-game delay and timing logic \\
\texttt{random}       & Python & Random ship placement in Battleship \\
\texttt{sys}          & Python & Clean process exit in Pong \\
\texttt{awk / grep / sort} & Bash & Text processing in leaderboard script \\
\bottomrule
\end{tabular}
\end{center}

\subsection{Pygame}
Pygame is a cross-platform Python library for building multimedia applications. In this project
it manages the display window, captures mouse and keyboard events, renders shapes and text,
and handles the game loop at 60 FPS.

\subsection{NumPy}
NumPy provides the \texttt{ndarray} data structure used to represent every game board. Its
vectorised Boolean operations allow efficient win detection without nested Python loops, keeping
game-state checks fast even on large boards (e.g.\ the $10{\times}10$ Tic-Tac-Toe grid).

\subsection{Matplotlib}
Matplotlib generates static charts from the game history CSV file: a pie chart of the
most-played games and a bar chart of the top five players by win count.

\subsection{Bash Utilities}
Standard POSIX tools (\texttt{awk}, \texttt{grep}, \texttt{cut}, \texttt{sort}, \texttt{uniq})
are used in \texttt{leaderboard.sh} to parse and aggregate the CSV history file without
requiring any additional external packages.

% -----------------------------------------------
\newpage
\section{Description of All Features Implemented}

\subsection{User Authentication System}
The authentication system, implemented entirely in Bash (\texttt{main.sh}), supports:
\begin{itemize}[leftmargin=*]
    \item \textbf{Registration} --- new users choose a username and password; the password is
          hashed with SHA-256 via \texttt{shasum -a 256} before being appended to \texttt{users.tsv}.
    \item \textbf{Login} --- returning users are authenticated by re-hashing their entered
          password and comparing it to the stored hash.
    \item \textbf{Two-player enforcement} --- the system ensures Player 2 must be a different
          account from Player 1, looping until a distinct user is authenticated.
    \item \textbf{Persistent storage} --- credentials are stored in a tab-separated file
          (\texttt{users.tsv}) that persists across sessions.
\end{itemize}

\subsection{Game Menu (game.py)}
After authentication, a Pygame-based menu allows the two authenticated players to select from
the available games. The menu launches the chosen game module and records the result (winner,
loser, game name) to \texttt{history.csv} upon completion.

\subsection{Five Playable Games}
\begin{enumerate}[leftmargin=*]
    \item \textbf{Tic-Tac-Toe} --- $10{\times}10$ board; five-in-a-row win condition; Player 1
          plays circles (O), Player 2 plays crosses (X).
    \item \textbf{Connect Four} --- $7{\times}7$ board; gravity-based token drop; four-in-a-row
          win in any direction.
    \item \textbf{Othello (Reversi)} --- $8{\times}8$ board; standard disc-flipping rules;
          winner determined by piece count at game end.
    \item \textbf{Battleship} --- $10{\times}7$ dual-grid board; randomised ship placement;
          players fire alternately; all enemy ships sunk wins.
    \item \textbf{Pong} --- real-time paddle-and-ball game; first to 5 points wins; ball angle
          varies with paddle contact point.
\end{enumerate}

\subsection{Common In-Game UI Elements}
All games share:
\begin{itemize}[leftmargin=*]
    \item Themed background images loaded from an \texttt{assets/} directory.
    \item A consistent retro font (\texttt{PressStart2P-Regular.ttf}).
    \item A \textbf{BACK} button (returns to menu without recording a result) and a
          \textbf{RESET} button (restarts the current game).
    \item A winner/draw pop-up overlay displayed 0.75\,s after the game ends,
          auto-dismissed after 2.5\,s.
    \item Mouse-hover colour change (white $\to$ yellow) on interactive buttons.
\end{itemize}

\subsection{Game History Recording}
Every completed game appends a row to \texttt{history.csv} containing the winner's name,
the loser's name, and the game title. This flat file acts as the single source of truth for
all leaderboard and analytics calculations.

\subsection{Leaderboard (leaderboard.sh)}
The leaderboard is a Bash script that:
\begin{itemize}[leftmargin=*]
    \item Reads \texttt{history.csv} and lists stats (wins, losses, win/loss ratio) per player,
          broken down by game.
    \item Displays the game passed as an argument first, followed by all other games found in
          the history.
    \item Supports three sort modes: by wins (\texttt{1}), by losses (\texttt{2}), or by
          win/loss ratio (\texttt{3}).
    \item Shows \texttt{UNDEFEATED} for players with no losses.
\end{itemize}

\subsection{Data Visualisations}
Two Matplotlib charts are generated from \texttt{history.csv}:
\begin{itemize}[leftmargin=*]
    \item \textbf{Most Played Games} --- a pie chart showing the percentage of matches played
          per game.
    \item \textbf{Top 5 Players} --- a bar chart ranking the five players with the highest win
          counts.
\end{itemize}

% -----------------------------------------------
\newpage
\section{Explanation of Code Logic}

\subsection{Generic Game Class (\texttt{game.py})}

All five games inherit from a single base class, \texttt{games\_class}, which encapsulates
shared state and behaviour:

\begin{lstlisting}[language=Python, caption={Core base class attributes}]
class games_class:
    def _init_(self, player1, player2, board=None):
        self.player1  = player1   # authenticated username
        self.player2  = player2
        self.board    = board     # numpy ndarray (game-specific shape)
        self.active   = 1         # whose turn (1 or 2)
        self.n_active = 2         # the other player
        self.winner   = 0         # 0 = no winner yet
        self.win_line = None      # used by Tic-Tac-Toe for line drawing

    def switch(self):
        self.active, self.n_active = self.n_active, self.active
\end{lstlisting}

Each subclass must implement three methods:
\begin{enumerate}
    \item \texttt{win\_check()} --- inspects the board and returns the winning player
          identifier or \texttt{0}.
    \item \texttt{show(screen, mouse\_pos)} --- renders the background, labels, and UI
          buttons for a single frame.
    \item \texttt{execution()} --- contains the main game loop; returns the winner
          (or \texttt{0} for draw/back).
\end{enumerate}

This structure eliminates code duplication and ensures that adding a new game requires
only implementing the three game-specific methods.

\subsection{Authentication System (\texttt{main.sh})}

\begin{lstlisting}[language=bash, caption={Password hashing and credential lookup}]
hash_password(){
    echo "$1" | shasum -a 256 | awk '{print $1}'
}

authenticate_user(){
    read -p "Username: " username
    stored=$(grep "^$username" "$users")
    if [[ -n $stored ]]; then
        read -sp "Password: " password
        hashed=$(hash_password "$password")
        stored_hash=$(echo "$stored" | cut -f2)
        [[ $hashed == $stored_hash ]] && echo "$username"
    fi
}
\end{lstlisting}

The function outputs the authenticated username to stdout, while all prompts and messages go
to stderr (\texttt{\&2}), allowing the calling code to capture only the clean username via
command substitution: \texttt{user1=\$(authenticate\_user 1)}.

\subsection{Tic-Tac-Toe (\texttt{tictactoe.py})}

The board is a $10{\times}10$ NumPy array. Win detection checks for five-in-a-row by
constructing Boolean truth matrices and combining them with bitwise AND. If we have five-in a row then that position in the final matrix will have true value and rest all will have false. Since this returns the position from where the five in a row sequence begins we can use it to extract the exact position where the win happens and a win line is drawn:

\begin{lstlisting}[language=Python, caption={Vectorised horizontal win check}]
t = self.board == self.active   # Boolean matrix
h = (t[:, 0:5] & t[:, 1:6] & t[:, 2:7]
       & t[:, 3:8] & t[:, 4:9])
if np.any(h):
    posn = np.where(h == 1)
    # record win_line endpoints for drawing
    return self.active
\end{lstlisting}

The same pattern repeats for vertical and both diagonal directions. A hover effect previews
each player's symbol in a dimmed colour before they click.

\subsection{Connect Four (\texttt{connect4.py})}

The $7{\times}7$ board uses values \texttt{1} (Player 1), \texttt{2} (Player 2), and
\texttt{-1} (winning tokens highlighted in orange). Gravity is handled by scanning a column
from the bottom row upward until an empty cell is found:

\begin{lstlisting}[language=Python, caption={Gravity: finding the lowest empty row}]
def available_row(self, column):
    for i in range(6, -1, -1):
        if not self.occ(i, column):
            return i
    return -1  # column is full
\end{lstlisting}

Win detection uses the same sliding Boolean window technique as Tic-Tac-Toe, but for windows
of length 4 in all four orientations. Winning cells are set to \texttt{-1} so they render
in a distinct highlight colour.

\subsection{Othello (\texttt{othello.py})}

Othello is the most algorithmically complex game. The core method is
\texttt{finding\_valid(box, check=False)}:

\begin{enumerate}
    \item It inspects the $3{\times}3$ neighbourhood around the clicked cell to find adjacent
          opponent pieces.
    \item For each direction containing an opponent piece, it walks along that direction until
          it hits a boundary, an empty cell (invalid), or the active player's own piece
          (valid --- flip all pieces in between).
    \item When \texttt{check=True}, it returns a Boolean immediately without modifying the
          board, used to pre-highlight valid moves and detect when no valid moves remain.
\end{enumerate}

When a player has no valid moves the turn is skipped; if neither player can move the game ends
and the winner is determined by piece count.

\subsection{Battleship (\texttt{battleship.py})}

Each player has an independent $10{\times}7$ NumPy board for ship positions and a separate
hits grid. Ships are placed randomly with collision detection:

\begin{lstlisting}[language=Python, caption={Ship placement with collision detection}]
def place_ships(self, board):
    for ship_size in [5, 4, 3, 3, 2]:
        placed = False
        while not placed:
            orientation = random.choice(["H", "V"])
            if orientation == "H":
                r = random.randint(0, 9)
                c = random.randint(0, 7 - ship_size)
                if np.all(board[r, c:c+ship_size] == 0):
                    board[r, c:c+ship_size] = 1
                    placed = True
\end{lstlisting}

The win condition checks whether all opponent ship cells have been hit:

\begin{lstlisting}[language=Python, caption={Win condition: all ships sunk}]
def win_check(self):
    if np.all((self.hits_p2 == 1) | (self.board_p2 == 0)):
        return self.player1
    elif np.all((self.hits_p1 == 1) | (self.board_p1 == 0)):
        return self.player2
    return 0
\end{lstlisting}

\subsection{Pong (\texttt{pong.py})}

Ball-paddle collisions compute a reflection angle proportional to how far the ball strikes
from the paddle centre. Ball speed increases by 5\% on every hit and is capped at 15
pixels/frame:

\begin{lstlisting}[language=Python, caption={Angle and speed modulation on paddle hit}]
if self.ball.colliderect(self.paddle1):
    self.ball_speed_x = abs(self.ball_speed_x)
    offset = ((self.ball.centery - self.paddle1.centery)
              / (self.paddle1.height // 2))
    self.ball_speed_y = offset * 6
    self.ball_speed_x *= 1.05
    self.ball_speed_y *= 1.05
    max_speed = 15
    self.ball_speed_x = max(-max_speed, min(max_speed, self.ball_speed_x))
    self.ball_speed_y = max(-max_speed, min(max_speed, self.ball_speed_y))
\end{lstlisting}

Player 1 uses \texttt{W}/\texttt{S} keys; Player 2 uses \texttt{UP}/\texttt{DOWN} arrow keys.
First to 5 points wins.

\subsection{Leaderboard (\texttt{leaderboard.sh})}

The script iterates over each unique game found in \texttt{history.csv}. For each game it
uses \texttt{awk} to count wins (rows where column 1 matches the player) and losses (rows
where column 2 matches). The ratio is computed with \texttt{awk} floating-point arithmetic.
Players with zero losses are marked \texttt{UNDEFEATED} with a sentinel ratio of
\texttt{999.00} so they always sort to the top.

% -----------------------------------------------
\newpage
\section{Hurdles Faced and How We Overcame Them}

\subsection{Cross-File Module Imports and Asset Paths}
Because each game file resides in a \texttt{games/} subdirectory while assets live in
\texttt{assets/}, hard-coded relative paths broke when the working directory changed. We
resolved this by computing absolute paths at runtime using \texttt{os.path.dirname(\\_file\\_)}
and \texttt{os.path.join}, making every import and image load independent of the invocation
directory.

\subsection{Othello Valid-Move Algorithm}
Correctly implementing Othello's flip logic along all eight directions and handling board edge
cases required several iterations. We refactored to use NumPy diagonal slicing
(\texttt{np.diagonal}) and \texttt{np.argmin} to locate the first non-opponent piece along
each ray, which simplified the code and eliminated boundary bugs.

\subsection{Bash and Python I/O Separation}
The \texttt{authenticate\_user} function needed to return a username while also printing
prompts to the terminal. Mixing stdout and stderr caused the captured output to include prompt
strings. The fix was to redirect all human-readable messages to \texttt{\&2} (stderr) so that
only the clean username was emitted on stdout and captured by command substitution.

\subsection{Carriage Returns in CSV}
On some systems \texttt{history.csv} contained Windows-style \texttt{\textbackslash r\textbackslash n}
line endings. The \texttt{awk} comparisons in \texttt{leaderboard.sh} failed silently because
the game name field included a trailing \texttt{\textbackslash r}. We added
\texttt{sub(/\textbackslash r\$/, "", \$3)} in every \texttt{awk} invocation to strip the
carriage return before any comparison.

\subsection{Battleship Turn Switching on Hits}
The original design switched turns after every click. In standard Battleship the shooter
continues on a successful hit. We added a conditional so \texttt{switch()} is called only when
the fired cell does not contain a ship, giving the correct ``keep going on a hit'' behaviour.

\subsection{Pygame Window Persistence Between Games}
Returning from one game to the menu while a Pygame display was still active caused the next
game to open a second window. We ensured each game's \texttt{execution()} method returns
cleanly and that the menu reclaims the single display via \texttt{pygame.display.set\_mode()},
avoiding duplicate windows.

% -----------------------------------------------
\newpage
\section{What We Would Do Differently With 10 More Hours}

\subsection{Save Good Games}
Add an option to record moves for Tic-tac-toe, Othello and Connect 4 and save these games to access later.

\subsection{Have an Option for Placing Ships in Battleship}
To stay true to the arcade theme we had the ships be randomised as players would have to look away as the other one places ships. We thought of adding an option to make this possible, upto the choice of the players.

\subsection{Animated UI and Sound Effects}
Adding Pygame sound effects (coin drops for Connect Four, explosion sounds for Battleship) and
smooth animations (token drop physics, disc flip for Othello) would significantly improve the
player experience with modest additional code.

\subsection{Responsive Layout}
All coordinates are currently hard-coded for a $900 \times 750$ pixel window. Abstracting
layout to proportional values would allow the games to scale to different screen sizes without
manual recalculation.

% -----------------------------------------------
\newpage
\section{Bibliography}

\begin{enumerate}[leftmargin=*]

    \item Pygame Community. \textit{Pygame Documentation}.
          \url{https://www.pygame.org/docs/}

    \item Harris, C.R. et al. (2020). Array programming with NumPy.
          \textit{Nature}, 585, 357--362.
          \url{https://doi.org/10.1038/s41586-020-2649-2}

    \item Hunter, J.D. (2007). Matplotlib: A 2D graphics environment.
          \textit{Computing in Science \& Engineering}, 9(3), 90--95.
          \url{https://doi.org/10.1109/MCSE.2007.55}

    \item Python Software Foundation. \textit{os --- Miscellaneous operating system interfaces}.
          \url{https://docs.python.org/3/library/os.html}

    \item Python Software Foundation. \textit{hashlib --- Secure hashes and message digests}.
          \url{https://docs.python.org/3/library/hashlib.html}

    \item GNU Project. \textit{Bash Reference Manual}.
          \url{https://www.gnu.org/software/bash/manual/bash.html}

    \item The Open Group. \textit{awk --- pattern scanning and processing language (POSIX)}.
          \url{https://pubs.opengroup.org/onlinepubs/9699919799/utilities/awk.html}

    \item World Othello Federation. \textit{Official Othello Rules}.
          \url{https://www.worldothello.org/about/about-othello/othello-rules/official-rules/english}

    \item Google Fonts. \textit{Press Start 2P}.
          \url{https://fonts.google.com/specimen/Press+Start+2P}

    \item NumPy Developers. \textit{numpy.diagonal reference}.
          \url{https://numpy.org/doc/stable/reference/generated/numpy.diagonal.html}

\end{enumerate}

\end{document}
